\SourceLeaf{248}%
\Needspace{7\baselineskip}
\begin{center}
[\textsc{Gerling} \emph{an} \textsc{Gauss}.\quad \emph{Marburg, 26. Februar 1844.}]
\end{center}

\{\mbox{.\ .\ .\ .\ .\ .\ .\ .\ .} Seitdem ich 1813 auf die symmetrischen Dreiecke und die Fläche
des sphärischen Dreiecks kam, ist mir immer das Desiderium geblieben, dass
mir eigentlich eine scharfe Definition oder ein scharfes Kriterium der
Symmetrie fehlte, im Gegensatz gegen die Congruenz. Seit vorigem Winter bin
ich nun darauf gekommen, die Sache auch im \textls{Vortrag}, denn früher kam ich
schon bei anderer Gelegenheit für mich darauf (1834), so darzustellen:

Jedes Gebilde in der Ebene beziehen wir auf \textls{zwei} Dimensionen (rechtwinklige Coordinaten), jedes räumliche auf
\textls{drei}. Denken wir nun die Coordinaten als dem Wechsel des Zeichens unterworfen,
so wird allgemein congruent, was in \textls{zwei} Dimensionen das Zeichen ändert;
symmetrisch aber, was in \textls{einer} oder in \textls{drei} Dimensionen das Zeichen ändert. --- Hieraus ergibt sich dann als Corollar,

1) dass bei Coordinaten aus einem Punkt durch Wechsel aller radii vectores
(sit venia verbo) Symmetrie entstehen muss bei körperlichen Constructionen,
Congruenz aber bei ebenen,

\Needspace{7\baselineskip}
2) dass die Symmetrie in die Congruenz übergeht, wenn das Gebilde durch eine
gerade Linie oder durch eine Ebene in zwei symmetrische Hälften theilbar
ist. ---

Sie würden mich sehr verbinden, wenn Sie mir gelegentlich Ihr Urtheil über
diese Darstellungsweise mittheilten.\}

\par\bigskip

\Needspace{7\baselineskip}
\begin{center}
\textsc{Gauss} an \textsc{Gerling}.\quad Göttingen, 8. April 1844.
\end{center}

\mbox{.\ .\ .\ .\ .\ .\ .\ .} Ihre Bemerkungen über Symmetrie und Congruenz sind vollkommen
treffend. Was noch zu desideriren wäre, ist der metaphysische Grund,
\textls{warum} es so ist (was bei Ihnen nur als eine wahrgenommene Thatsache auftritt)
und damit auch die Erweiterung auf eine Geometrie von mehr als 3
Dimensionen, wofür wir menschliche Wesen keine Anschauung haben, die aber
